\section{Adequacy and the Cut Boundary of the Exact Equality LK}
\label{sec:p07-equality-lk-metatheory}

The equality extension on \hyperref[sheet:lk-eq]{Sheet~\ref*{sheet:lk-eq}}
is sound and complete when Cut is available. Cut is nevertheless not
admissible in its Cut/Mix-free fragment. The distinction is caused by the
particular orientation and placement of its equality rules, not by a failure
of identity semantics. This section supplies the complete translations and
the explicit negative instance for that existing surface. It changes neither
the sheet nor the inherited equality presentation in
\cref{sec:fol-deduction-02} \cite{Enderton2001,Mendelson1997}.

\begin{hbdefinition}{def:p07-eqlk-domain}{The literal equality calculus and its validity}
Write \(\mathsf{LK}_{=}\) for exactly \texttt{sheet:lk-eq}.
Formulas have primitive \(\bot,\neg,\land,\lor,\to,\forall,\exists\)
and identity atoms. Both contexts are finite multisets. Identity and
falsum initials are contexted. Reflexivity has the empty-context initial
\(\Rightarrow t=t\). The only primitive equality replacement is the
left rule \ref{rule:lk-eq-repl-l}; its direction, retained equality premise
and untouched succedent are part of the definition. Cut and positive-
multiplicity Mix are the existing registered rules, not admissibility claims.
A \emph{Cut/Mix-free derivation} uses neither of those two rules.

Every term instance is the total capture-avoiding substitution of
\cref{def:p02-fol-capture-avoiding-substitution}. The quantifier rules
retain their exact eigenvariable freshness conditions. No implicit
formula-conversion, right equality rule, reflexivity-deletion rule, or
opposite-oriented left replacement rule is added.

The semantic structures are the nonempty structures of
\cref{sec:fol-semantics-01}; equality is actual identity and functions
are total. A sequent is valid when, at every structure and assignment,
truth of all antecedent formulas entails truth of at least one succedent
formula. Thus an empty succedent requires that the antecedent be
unsatisfiable, and an empty antecedent requires validity of the succedent
disjunction. For \(\Delta\), write \(D_\Delta\) for the finite-support
disjunction fixed in \cref{def:p07-sequent-nd-translation}, with
\(D_\varnothing=\bot\).
\index{equality!sequent calculus}\index{cut admissibility!failure}
\end{hbdefinition}

\begin{hblemma}{thm:p07-eqlk-alpha}{proved=annex:p7.eqlk.alpha}{Literal alpha-conversion kernels with equality}
If \(A\equiv_\alpha B\), the exact equality sheet has Cut/Mix-free
derivations of \(A\Rightarrow B\) and \(B\Rightarrow A\).
Consequently a cut-admitting derivation can convert a chosen antecedent
or succedent occurrence to an alpha-equivalent literal formula by a Cut
with the appropriate kernel. No alpha rule is presumed on the sheet.
\end{hblemma}

\begin{hblemma}{thm:p07-eqlk-replacement}{proved=annex:p7.eqlk.replacement}{Symmetry and the two substitution directions}
For arbitrary terms \(r,t\), formula \(A\), and variable \(x\),
the exact sheet derives
\[
 r=t\Rightarrow t=r,
 \qquad
 r=t,A[t/x]\Rightarrow A[r/x],
 \qquad
 r=t,A[r/x]\Rightarrow A[t/x].
\]
The middle kernel is Cut/Mix-free; the symmetry construction and the
last kernel explicitly use Cut. The opposite-oriented left replacement
of a proof premise is derived with those Cuts. Terms may contain \(x\),
and \(A\) may have zero or multiple free occurrences of \(x\).
\end{hblemma}

\begin{hbtheorem}{thm:p07-eqlk-congruence}{proved=annex:p7.eqlk.congruence}{Transitivity and function/relation congruence}
The exact sheet derives \(r=t,t=u\Rightarrow r=u\).
For every finite arity \(n\), terms \(r_i,t_i\), function symbol \(f\)
of arity \(n\), and relation symbol \(R\) of arity \(n\), it derives
\[
 r_1=t_1,\ldots,r_n=t_n
 \Rightarrow f(r_1,\ldots,r_n)=f(t_1,\ldots,t_n),
\]
\[
 r_1=t_1,\ldots,r_n=t_n,R(r_1,\ldots,r_n)
 \Rightarrow R(t_1,\ldots,t_n),
\]
and the relation kernel with the two argument tuples interchanged.
For \(n=0\), the function case is reflexivity of the constant and the
relation case is identity of the nullary atom. These are cut-admitting
constructions, not a blanket Cut-free claim.
\end{hbtheorem}

\begin{hbtheorem}{thm:p07-eqlk-soundness}{proved=annex:p7.eqlk.soundness}{Soundness of every literal equality-LK rule}
Every \(\mathsf{LK}_{=}\)-derivable sequent is valid in the nonempty
identity structures and assignment semantics of
\cref{def:p07-eqlk-domain}. This covers arbitrary finite multiset contexts,
empty contexts, all total substitution instances, all four quantifier
rules, both contractions, Cut and the exact positive-multiplicity Mix.
\end{hbtheorem}

\begin{hbtheorem}{thm:p07-eqlk-nd}{proved=annex:p7.eqlk.nd}{Rule-by-rule equivalence with identity natural deduction}
For every finite multiset sequent \(\Gamma\Rightarrow\Delta\),
\(\mathsf{LK}_{=}\) derives it exactly when
\hyperref[sheet:nd-fol]{the exact classical identity natural-deduction sheet}
derives \(D_\Delta\) with open assumptions in
\(\operatorname{supp}(\Gamma)\). The constructions preserve the total
substitution and eigenvariable conditions and respect the natural-
deduction sheet's nonvacuous discharge requirements. They establish
derivability equivalence, not equality of proof trees or preservation
of Cut-freeness.
\end{hbtheorem}

\begin{hbtheorem}{thm:p07-eqlk-completeness}{proved=annex:p7.eqlk.completeness}{Finite-sequent identity completeness for the unchanged sheet}
For every signature and finite sequent, validity in the structures and
assignments of \cref{def:p07-eqlk-domain} is equivalent to derivability
on the exact existing equality sheet, with Cut admitted. No ambient
Axiom of Choice is needed for this finite-sequent statement: reduce
the semantic question to the finite signature occurring in the sequent
and use the proved countable-signature identity completeness theorem.
This is neither a Cut-free completeness result nor a claim that
arbitrary theories over arbitrary signatures have the same choice boundary.
\end{hbtheorem}

\subsection{The negative Cut result}

\begin{hbtheorem}{thm:p07-eqlk-cut-failure}{proved=annex:p7.eqlk.cut-failure}{Cut and Mix are not admissible for this equality surface}
In a signature containing two syntactically distinct constants \(a,b\),
the valid sequent
\[
                         a=b\Rightarrow b=a
\]
has a derivation with one Cut between two Cut/Mix-free premises, but
has no Cut/Mix-free derivation on \texttt{sheet:lk-eq}.
The same construction and obstruction use two distinct free variables
instead, so the negative instance does not require a signature extension.
Consequently deleting Cut and Mix changes the set of derivable sequents;
neither Cut nor Mix is admissible in the resulting fragment, and
semantic completeness fails for that fragment.
\end{hbtheorem}

\begin{hbwarning}{def:p07-eqlk-cut-boundary}{Adequacy is not analytic completeness}
The symmetry result above is proved \emph{with} Cut. It cannot be used
as a Cut-free permutation subroutine: the displayed counterexample proves
that such a subroutine does not exist for this sheet. Contraction makes
the registered positive Mix and Cut interderivable when they are allowed;
it does not make either admissible when both are removed.
The defined-negation equality-free Cut theorem
\cref{thm:p07-classical-cut-elimination} and the equality-free
canonical-LK bridge concern other surfaces. Adding a right substitution
rule, reflexivity deletion, or a bidirectional left rule would define a
distinct calculus requiring its own adequacy and Cut analysis. None is
silently added here.
\end{hbwarning}
